\documentclass[12pt,a4paper]{report}

% Packages
\usepackage[utf8]{inputenc}
\usepackage{graphicx}
\usepackage{amsmath, amssymb}
\usepackage[a4paper, top=0.5cm, bottom=2.5cm, left=2.5cm, right=2.5cm]{geometry}
\usepackage{setspace}
\usepackage{hyperref}
\usepackage{fancyhdr}
\usepackage{titlesec}
\usepackage{relsize}
\usepackage{enumitem}
\usepackage[most]{tcolorbox}


% Page setup
% \geometry{margin=1in}
\onehalfspacing
\setlength{\parskip}{0.8em}
\setlength{\parindent}{0pt}

% Header & Footer
\pagestyle{fancy}
\fancyhf{}
\fancyhead[L]{\vspace{1.5cm}\raisebox{-40pt}[0pt][0pt]{Pràctica 2 - Aprenentatge per reforç}} % \fancyhead[L]{Your Report Title}
\fancyhead[R]{\vspace{1.5cm}\raisebox{-40pt}[0pt][0pt]{\thepage}} % \fancyhead[R]{\thepage}

\setlength{\headheight}{10pt} %\setlength{\headheight}{10pt}
\setlength{\headsep}{70pt}

% Title setup
\titleformat{\chapter}[display]
  {\normalfont\bfseries\Huge}
  {\vspace{-2cm}}{20pt}{\Huge}

% Title Page
\title{
  \vspace{0cm}
  \textbf{Memòria de la Pràctica 2}\\
  {\Large \textbf{Aprenentatge per reforç}}
  \vspace{5cm}
}
\author{Serafí Nebot Ginard \\  Jaume Galmés Ramis \\
\vspace{7cm}
\\ 21722 - Intel·ligència Artificial \\ Grau d’Enginyeria Informàtica \\ Curs 2025/2026
}

\fancypagestyle{plain}{
  \fancyhf{} % neteja
  \fancyfoot[C]{\thepage} % número de pàgina al centre
  \renewcommand{\footrulewidth}{0pt}
  \setlength{\footskip}{1.5cm} % ajusta la distància des del final de la pàgina
}

\begin{document}

\date{}  % No mostrar cap data
\maketitle
\thispagestyle{empty}
\newpage
\titlespacing*{\chapter}{0pt}{0pt}{20pt}  

\pagenumbering{roman}
\renewcommand{\contentsname}{Índex}
\tableofcontents
\newpage

\pagenumbering{arabic}

% ==========================
\chapter{Introducció}
\label{ch:intro}

\dots

Investigant diferents algorismes de RL hem trobat el Average Q-learning, que és una extensió del Q-learning que hem vist a classe.
Ens ha semblat interessant i l'hem inclòs al projecte.


% ==========================
\chapter{Disseny algorismes}
\label{ch:structure}

\section{Genètic}

\textbf{Individu}

Un individu representa una política $\pi_i(s_t) \in A$ on a cada possible estat se li assigna una acció.
La popblació d'individus conté, per a cada generació, $n$ individus.
Per a la primera generació, els gens de tots els individus (cromosomes) s'inicialitzen de forma completament aleatòria.

\textbf{Funció \textit{fitness}}

La funció fitness és la tasa d'èxit de l'individu. Al ser un entorn estocàstic s'evalua la política durant 100 episodis i es calcula el promig de vegades que s'ha arribat al destí.

Pel cas particular del \textit{Frozen Lake} totes les recompenses són 0 excepte quan s'arriba al destí, per tant, la funció \textit{fitness} és equivalent al promig de recompenses totals obtingudes al llarg de tots els episodis.

\[
  \mathit{fitness}({\pi_i}) = \frac{1}{100} \sum\limits_{j = 1}^{100} \sum\limits_{t = 1}^{T} R_t^j
\]

On $R_t^j$ és la recompensa del període $t$ de l'episodi $j$.

\textbf{Elitisme o \textit{Culling}}

Abans de generar la nova generació es seleccionen els $l = \lceil c \cdot n \rceil$ millors individus que passen directament a la nova generació (la resta es descarten), on $c \in [0, 1]$ és la tasa de \textit{culling}.
Una $c > 0$ ens garanteix que no es perdin els millors individus amb el pas de les generacions.

\textbf{Selecció}

De tots els individus de cada generació, $n$ en total, només es seleccionen els $k = \lceil \beta \cdot n \rceil$ millors, on $\beta \in (0, 1]$ és el paràmetre de pressió de selecció.
D'aquesta forma, $\beta$, ens permet escollir quin percentatge de la millor part de la població volem emprar per a generar la pròcima generació.

\textbf{Mutació}

Després de generar cada individu hi ha una petita probabilitat $m$ de mutar-lo de forma aleatòria.
Per fer-ho, per a cada gen de l'individu es genera un nombre aleatori $p \in [0, 1]$ i es genera un gen aleatori si $p < m$ i l'original si $p \geq m$.

\[
  \begin{cases}
    \textrm{original } & \textrm{ si } p \geq m \\
    \textrm{aleatori } & \textrm{ si } p < m \\
  \end{cases}
\]

\section{Programació dinámica}

\section{Monte Carlo}

\section{SARSA}

\section{Q-learning}

\pagebreak

% ==========================
\chapter{Estudi dels resultats}
\label{ch:results}

A continuació es mostra com diferents variacions dels paràmetres de cada un dels models afecta als resultats.

Per a cada configuració dels paràmetres, cada un dels algorismes primer passa per a la fase d'entranament de
$n$ episodis i després s'evalua la política resultant amb 1000 episodis per a contrarestar el factor estocàstic de l'entorn.

Per a mantenir el document breu, les úniques métriques que comparam són el nombre promig de passes (\textit{steps}) que es tarda per a cada episodi
i el promig de vegades que s'ha arribat al destí, anomenat \textit{sr} (\textit{Success Rate}).

Per a \textit{Monte Carlo}, \textit{SARSA} i \textit{Q-learning} tenim els següent paràmetres:
\begin{itemize}
  \item $\mathit{dr} = 0.99$: \textit{Discount Rate} o \textit{gamma} als apunts, és el factor de descompte per
    a la recompensa esperada per a cada passa del mateix episodi.
\end{itemize}

\begin{minipage}[t]{0.4\textwidth}
  \centering
  \begin{tabular}{| c | c c | c c |}
    \hline
      & \multicolumn{2}{|c|}{SARSA} & \multicolumn{2}{|c|}{Q-learning} \\
    \hline
    \textit{lr} & steps & sr & steps & sr \\
    \hline
    0.10 & 44.607 & 0.74  & 44.711 & 0.721 \\
    0.20 & 44.706 & 0.644 & 45.731 & 0.745 \\
    0.30 & 44.441 & 0.72  & 45.358 & 0.732 \\
    0.40 & 46.921 & 0.638 & 45.021 & 0.725 \\
    0.50 & 46.462 & 0.734 & 43.33  & 0.734 \\
    0.60 & 28.884 & 0.244 & 44.917 & 0.735 \\
    0.70 & 62.305 & 0.561 & 42.406 & 0.756 \\
    0.80 & 61.438 & 0.573 & 44.273 & 0.758 \\
    0.90 & 8.856  & 0.102 & 45.38  & 0.744 \\
    1.00 & 7.832  & 0.024 & 7.837  & 0.022 \\
    \hline
  \end{tabular}
\end{minipage}
\hfill
\begin{minipage}[t]{0.5\textwidth}
  Podem veure que SARSA és molt més sensible al \textit{Learning Rate} que Q-learning,
  ja que aquest darrer es manté igual al llarg de totes les combinacions.

  Per a valors molt grans de \textit{lr} ambdós tenen resultats molt pobres, sobretot quan
  son molt propers a 1. Això té sentit perque significa que per a cada actualització de la política s'està
  corregint un 100\% de l'error, es a dir, s'estàn fent salts molt grans i no s'arriba a la solució òptima
  (s'arriba a un máxim local enlloc d'un màxim global).
\end{minipage}

\pagebreak

\begin{table}
  \centering
  \begin{tabular}{| c | c c | c c | c c |}
    \hline
      & \multicolumn{2}{|c|}{Monte Carlo} & \multicolumn{2}{|c|}{SARSA} & \multicolumn{2}{|c|}{Q-learning} \\
    \hline
    \textit{dr} & steps & sr & steps & sr & steps & sr \\
    \hline
    0.10 & 13.4970 & 0.1060 & 9.2340  & 0.0620 & 9.2420  & 0.1120 \\
    0.20 & 7.8380  & 0.0470 & 8.5430  & 0.1010 & 7.3150  & 0.0470 \\
    0.30 & 7.5250  & 0.0650 & 16.0000 & 0.1760 & 6.2980  & 0.0600 \\
    0.40 & 10.9760 & 0.0860 & 20.9020 & 0.1790 & 14.2070 & 0.1560 \\
    0.50 & 35.9780 & 0.6950 & 14.5330 & 0.1570 & 21.9130 & 0.2800 \\
    0.60 & 10.0380 & 0.0930 & 33.7380 & 0.5140 & 22.0330 & 0.2630 \\
    0.70 & 9.1360  & 0.0890 & 16.2900 & 0.2180 & 30.5880 & 0.4590 \\
    0.80 & 33.2870 & 0.4890 & 26.8600 & 0.3900 & 27.9760 & 0.4460 \\
    0.90 & 30.6570 & 0.5020 & 41.5030 & 0.6280 & 41.8290 & 0.7180 \\
    1.00 & 11.2490 & 0.1090 & 44.7220 & 0.7600 & 100.000 & 0.0000 \\
    \hline
  \end{tabular}
\end{table}

\begin{table}
  \centering
  \begin{tabular}{| c | c c | c c | c c |}
    \hline
      & \multicolumn{2}{|c|}{Monte Carlo} & \multicolumn{2}{|c|}{SARSA} & \multicolumn{2}{|c|}{Q-learning} \\
    \hline
    \textit{er} & steps & sr & steps & sr & steps & sr \\
    \hline
    0.10 & 24.7740  & 0.2420 & 43.8000 & 0.7430  & 46.1800 & 0.7360 \\
    0.20 & 27.6060  & 0.1950 & 45.8280 & 0.7440  & 43.3920 & 0.7630 \\
    0.30 & 38.0640  & 0.6740 & 43.5510 & 0.7550  & 45.3150 & 0.7380 \\
    0.40 & 30.1160  & 0.5200 & 45.4450 & 0.7440  & 45.6250 & 0.7290 \\
    0.50 & 40.1320  & 0.7390 & 43.4110 & 0.7570  & 45.3650 & 0.7710 \\
    0.60 & 30.0940  & 0.2350 & 43.7320 & 0.7640  & 45.3130 & 0.6970 \\
    0.70 & 40.6120  & 0.7150 & 43.8760 & 0.7380  & 44.4060 & 0.7330 \\
    0.80 & 40.6070  & 0.7340 & 44.9870 & 0.7380  & 44.1360 & 0.7530 \\
    0.90 & 41.1670  & 0.7310 & 42.6170 & 0.7520  & 45.2820 & 0.7320 \\
    1.00 & 41.2560  & 0.7030 & 45.6850 & 0.7440  & 45.0200 & 0.7330 \\
    \hline
  \end{tabular}
\end{table}

\begin{table}
  \centering
  \begin{tabular}{| c | c c | c c | c c |}
    \hline
      & \multicolumn{2}{|c|}{Monte Carlo} & \multicolumn{2}{|c|}{SARSA} & \multicolumn{2}{|c|}{Q-learning} \\
    \hline
    \textit{er min} & steps & sr & steps & sr & steps & sr \\
    \hline
    0.10 & 11.1730 & 0.1660 & 40.7290 & 0.7200 & 45.1450 & 0.7380 \\
    0.20 & 40.7990 & 0.7040 & 39.1340 & 0.7460 & 45.5910 & 0.7680 \\
    0.30 & 40.7220 & 0.7460 & 41.3910 & 0.7160 & 44.2060 & 0.7570 \\
    0.40 & 30.2160 & 0.5120 & 41.9450 & 0.7300 & 43.6270 & 0.7450 \\
    0.50 & 39.6190 & 0.6070 & 39.7820 & 0.6890 & 43.4790 & 0.7480 \\
    0.60 & 40.5510 & 0.7290 & 42.2630 & 0.7120 & 45.9040 & 0.7170 \\
    0.70 & 26.8700 & 0.3990 & 23.1390 & 0.2830 & 44.0580 & 0.7150 \\
    0.80 & 30.4010 & 0.5150 & 39.8010 & 0.7070 & 44.8900 & 0.7240 \\
    0.90 & 18.4920 & 0.1910 & 36.0990 & 0.6080 & 44.2430 & 0.7470 \\
    1.00 & 12.9100 & 0.1410 & 24.2430 & 0.3300 & 45.1970 & 0.7510 \\
    \hline
  \end{tabular}
\end{table}

\begin{table}
  \centering
  \begin{tabular}{| c | c c | c c | c c |}
    \hline
      & \multicolumn{2}{|c|}{Monte Carlo} & \multicolumn{2}{|c|}{SARSA} & \multicolumn{2}{|c|}{Q-learning} \\
    \hline
    \textit{er decay} & steps & sr & steps & sr & steps & sr \\
    \hline
    0.10 & 8.0950  & 0.0810 &         &        &         &        \\
    0.20 & 24.4520 & 0.2640 &         &        &         &        \\
    0.30 & 27.6170 & 0.4670 &         &        &         &        \\
    0.40 & 30.0460 & 0.2070 &         &        &         &        \\
    0.50 & 10.5710 & 0.0960 &         &        &         &        \\
    0.60 & 24.0870 & 0.1480 &         &        &         &        \\
    0.70 & 13.6050 & 0.1610 &         &        &         &        \\
    0.80 & 38.2980 & 0.6750 &         &        &         &        \\
    0.90 & 9.6760  & 0.0970 &         &        &         &        \\
    1.00 & 39.8340 & 0.6960 &         &        &         &        \\
    \hline
  \end{tabular}
\end{table}

% ==========================
\chapter{Conclusió}
\label{ch:conclusion}

\cleardoublepage
\thispagestyle{empty} % sense número de pàgina

\begin{center}
\vspace*{5cm} % espai superior

{\LARGE \textbf{Pràctica Bloc I}}\\[1em]
{\Large Intel·ligència Artificial}\\[4cm]

{\large Aprenentatge per reforç}\\[5cm]

{\large %\textbf{Alumnes:}}\\
Serafí Nebot Ginard\\
Jaume Galmés Ramis\\[4cm]}

{\large Grau d'Enginyeria Informàtica\\Curs 2025/2026}

\end{center}

\end{document}